\section*{Bab \ref{ch:g4:magnets-and-compass} --- Magnet dan Kompas}

\begin{solution}{exo:g4:magnets-and-compass:1}
\omterm{def:g4:magnets-and-compass:poles}{Kutub} adalah kedua ujung \omterm{def:g1:magnets:magnet}{magnet} yang kuat: \omterm{def:g4:magnets-and-compass:poles}{kutub utara} (U, sering
merah) dan \omterm{def:g4:magnets-and-compass:poles}{kutub selatan} (S). Penjepit kertas berkerumun di \omterm{def:g4:magnets-and-compass:poles}{kutubnya}
--- di kedua ujungnya --- bukan di tengah.
\end{solution}

\begin{solution}{exo:g4:magnets-and-compass:2}
\omterm{def:g4:magnets-and-compass:poles}{Kutub} yang berlawanan saling menarik; \omterm{def:g4:magnets-and-compass:poles}{kutub} yang sejenis saling
menolak. Tempelan berarti \omterm{def:g4:magnets-and-compass:poles}{kutub} yang berlawanan tadi berhadapan: U
terhadap S.
\end{solution}

\begin{solution}{exo:g4:magnets-and-compass:3}
U terhadap U: keduanya terdorong menjauh (\omterm{def:g4:magnets-and-compass:poles}{kutub} sejenis saling
menolak). U terhadap S: keduanya saling menarik lalu menempel.
\end{solution}

\begin{solution}{exo:g4:magnets-and-compass:4}
Karena ada \omterm{def:g1:magnets:magnet}{magnet} raksasa yang menarik \omterm{def:g4:magnets-and-compass:poles}{kutubnya} sampai sejajar: Bumi
sendirilah \omterm{def:g1:magnets:magnet}{magnet} itu, dengan \omterm{def:g4:magnets-and-compass:poles}{kutub} \omterm{def:g1:magnets:magnet}{magnetnya} di dekat puncak dan
dasar bola dunia.
\end{solution}

\begin{solution}{exo:g4:magnets-and-compass:5}
Utara, Timur, Selatan, Barat. Menghadap utara, Timur ada di sebelah kananmu.
\end{solution}

\begin{solution}{exo:g4:magnets-and-compass:6}
Jarumnya adalah \omterm{def:g1:magnets:magnet}{magnet} mungil dan menurut kepada tarikan kuat yang
\emph{paling dekat}. \omterm{def:g1:magnets:magnet}{Magnet} atau bongkahan besi besar di dekatnya
mengalahkan tarikan \omterm{def:g1:magnets:magnet}{magnet} Bumi yang jauh, dan jarumnya pun berbohong
tentang utara.
\end{solution}

\begin{solution}{exo:g4:magnets-and-compass:7}
Gosokkan jarumnya dengan satu \omterm{def:g4:magnets-and-compass:poles}{kutub} \omterm{def:g1:magnets:magnet}{magnet}, dua puluh kali ke arah
yang sama, sambil mengangkat \omterm{def:g1:magnets:magnet}{magnetnya} di antara gosokan --- jarum itu
menjadi \omterm{def:g1:magnets:magnet}{magnet} kecil. Letakkan di atas irisan gabus yang mengapung di
baskom berisi air: gabus yang mengapung itu poros yang terbuat dari
air, dan jarumnya berayun sampai menunjuk utara--selatan.
\end{solution}

\begin{solution}{exo:g4:magnets-and-compass:8}
\omterm{def:g1:light-and-shadows:shadow}{Bayang-bayang} itu menunjuk ke utara. \omterm{def:g4:magnets-and-compass:compass}{Kompas} yang jarumnya sependapat
dengan \omterm{def:g1:light-and-shadows:shadow}{bayang-bayang} tengah hari sedang mengatakan kebenaran --- dan
pada hari yang cerah \omterm{def:g1:light-and-shadows:shadow}{bayang-bayangnya} saja sudah dapat menjadi \omterm{def:g4:magnets-and-compass:compass}{kompas}
kasar.
\end{solution}

\begin{solution}{exo:g4:magnets-and-compass:9}
Ia bekerja dalam kabut, di bawah awan, pada malam tanpa bulan, dan
jauh dari pantai mana pun --- jarumnya tidak meminta apa pun dari
langit. Pelaut akhirnya dapat menjaga haluan menyeberangi samudra
terbuka dalam cuaca apa pun.
\end{solution}

\begin{solution}{exo:g4:magnets-and-compass:10}
Kepala ikat pinggang besi temannya mengalahkan \omterm{def:g1:magnets:magnet}{magnet} Bumi pada jarak
dekat: jarumnya menunjuk ke kepala ikat pinggang itu, bukan ke utara.
Kalau ia memercayainya, ia akan melenceng --- ``utara'' yang palsu itu
terletak ke arah kirinya, sehingga ia akan membelokkan jalannya ke
kiri, ke barat. Perbaikannya: bacalah kompasnya jauh dari ikat
pinggang itu (beberapa langkah terpisah), baru berjalan.
\end{solution}

\begin{solution}{exo:g4:magnets-and-compass:11}
Jarum pendek yang digosok itu punya \emph{dua} \omterm{def:g4:magnets-and-compass:poles}{kutub} --- kedua
ujungnya berperilaku sebagai U dan S (apungkan dua jarum semacam itu
dan keduanya menempel ujung ke ujung, ujung yang berlawanan saling
menarik). Jadi memendekkan \omterm{def:g1:magnets:magnet}{magnet} tidak menghilangkan satu \omterm{def:g4:magnets-and-compass:poles}{kutub}, dan
memotongnya menjadi dua hanya menghasilkan dua \omterm{def:g1:magnets:magnet}{magnet} pendek yang
lengkap: berburu \omterm{def:g4:magnets-and-compass:poles}{kutub} tunggal dengan cara memotong tidak pernah
berakhir --- setiap potongan menyerahkan kedua \omterm{def:g4:magnets-and-compass:poles}{kutubnya} lagi kepadamu.
\end{solution}
